% TOP_PROBLEM: {"id": 2302022, "title": "Research Problems in Function Theory — Problem 2.22", "category_id": 8, "category": {"id": 8, "name": "analysis", "display_name": "Analysis", "description": "Limits, continuity, calculus, and function theory.", "slug": "analysis", "order_index": 8, "created_at": "2026-07-31T15:26:25.671Z"}, "status": "open"}
% FIRST_POSTED: null
% ENTRY_KIND: statement_only
% Imported from: batch11.tex; UnsolvedMath ID 2302022
% Compile twice: lualatex 2302022.tex
\documentclass[11pt]{article}
\usepackage{amsmath,amssymb,mathtools}
\usepackage{fontspec}
\usepackage{unicode-math}
\setmainfont{Latin Modern Roman}
\setmathfont{Latin Modern Math}
\usepackage{xurl}
\usepackage[unicode,colorlinks=true,linkcolor=blue,urlcolor=blue,pdftitle={UnsolvedMath Problem 2302022}]{hyperref}
\newcommand{\sref}[2]{\hyperlink{source:#1:#2}{[S#2]}}
\newcommand{\cataloguescope}[1]{\paragraph{Mathematical statement.}#1}
\begin{document}
% UnsolvedMath ID 2302022; source catalogue code AMR-022-2022
\setcounter{section}{2302021}
\section{Research Problems in Function Theory — Problem 2.22}\label{2302022}
{\small\sffamily UnsolvedMath ID 2302022 · Analysis}
\subsection{Definitions and mathematical statement}

\paragraph{Shared notation.}$\mathbb N=\{1,2,\ldots\}$, and $\mathbb Z,\mathbb Q,\mathbb R,\mathbb C$ denote the integers, rationals, reals and complex numbers. Any other conventions are specified locally.


Let $f:\mathbb C\to\mathbb C$ be entire and set $f^{\circ1}=f$, $f^{\circ(n+1)}=f\circ f^{\circ n}$.
A point $z_0$ has exact period $n$ if $f^{\circ n}(z_0)=z_0$ and $f^{\circ k}(z_0)\ne z_0$ for $1\le k<n$.
The source calls such a point a ``fixed point of exact order $n$''.
Let $J(f)$ be the set of points with no neighbourhood on which the family $(f^{\circ n})_{n\ge1}$ is normal, where normality means that every sequence has a subsequence converging locally uniformly in the spherical metric. The source denotes this nonnormality set by $\mathcal F(f)$.

\paragraph{Statement recorded in the supplied source.}
Is there an entire function $f$ for which $J(f)=\mathbb C$? In particular, is $J(e^z)=\mathbb C$, equivalently, are the periodic points of $z\mapsto e^z$ dense in $\mathbb C$? \sref{2302022}{2}
Update 2.22 records Baker's example and Misiurewicz's proof for $f(z)=e^z$, answering both questions affirmatively. \sref{2302022}{2}
\subsection{Short English statement}
Can the iterates of an entire function fail to form a normal family everywhere, and does the exponential function already have this property?
\subsection{Sources}
\begin{description}
\item[\hypertarget{source:2302022:1}{S1}] ULAM.AI and the UnsolvedMath Contributors, \emph{UnsolvedMath: A Curated Collection of Open Mathematics Problems}, record 2302022 (AMR-022-2022). CC BY 4.0. \url{https://huggingface.co/datasets/ulamai/UnsolvedMath}.
\item[\hypertarget{source:2302022:2}{S2}] W. K. Hayman and E. F. Lingham, \emph{Research Problems in Function Theory} (2018), Chapter 2, Problem 2.22 and Update 2.22. \url{https://arxiv.org/abs/1809.07200}.
\end{description}
\label{end:2302022}
\end{document}
